% =====================================================================
\chapter{Segundo instrumento: FinBETO, el termómetro emocional}
\chaptermark{FinBETO, el termómetro emocional}
\label{cap:finbeto}

El corpus dice qué se publicó; falta saber con qué tono. Para eso se construyó FinBETO, un
modelo de lenguaje que lee cada oración de una nota y la clasifica como positiva, negativa,
neutra o mixta. Agregadas por rueda, esas clasificaciones dan una medida de la temperatura
emocional de la prensa económica: cuán optimistas o pesimistas fueron las noticias de cada día.

% ---------------------------------------------------------------------
\section{De BERT a FinBETO}

BERT \citep{devlin2019bert} es un modelo de lenguaje preentrenado con la tarea de adivinar
palabras ocultas en un texto (\textit{masked language modeling}, MLM). BETO es su versión
entrenada en español \citep{canete2020beto}. Para clasificar sentimiento se le agrega una capa
de salida y se lo ajusta con oraciones etiquetadas (\textit{fine-tuning}), como hizo FinBERT en
inglés \citep{araci2019finbert}. Antes del ajuste conviene seguir el preentrenamiento con texto
del propio dominio, porque mejora el desempeño en la tarea final \citep{gururangan2020dapt}.
FinBETO es el resultado de los dos pasos sobre la prensa económica argentina.

\textbf{Clases.} Se trabajó con cuatro: positiva (POS), negativa (NEG), neutra (NEU) y mixta
(MIXTO). Una parte de las oraciones de la prensa económica combina señales de ambos signos —por
ejemplo, un dato de actividad que mejora mientras la inflación se acelera— y forzarlas a una
sola clase habría sido arbitrario.

% ---------------------------------------------------------------------
\section{Etiquetado y entrenamiento}

\textbf{Etiquetado con Gemma.} Para el ajuste fino hacen falta oraciones etiquetadas, y no
existía un conjunto público de prensa económica argentina. Se usó un modelo de lenguaje como
etiquetador, una práctica que puede igualar o superar a anotadores contratados
\citep{gilardi2023chatgpt} y que también se evaluó en textos financieros
\citep{aguda2024llmannot}. El etiquetador fue Gemma \citep{gemma2026gemma4}
(\texttt{gemma4:26b}, temperatura 0,1), ejecutado localmente con Ollama, en dos pasos por
noticia. Primero leyó el título, la bajada y los primeros 3.000 caracteres, y produjo un
análisis de contexto (tema central, encuadre, actores y tono general). Después etiquetó cada
oración de al menos 30 caracteres, con ese contexto como referencia y la instrucción de que el
sentimiento surgiera de lo que la oración expresa por sí misma. Se tomó un artículo de cada una
de las nueve fuentes por lote, hasta 408 noticias (45 o 46 por fuente) y 11.605 oraciones
etiquetadas: 46,4\% NEU, 23,8\% POS, 20,8\% NEG y 8,9\% MIXTO.

\textbf{Preentrenamiento continuado.} Se hizo con la tarea MLM (15\% de los tokens
enmascarados) sobre el texto de los 14.667 artículos del corpus antes de su depuración final:
3 épocas, segmentos de 128 tokens y tasa de aprendizaje $2\cdot10^{-5}$. La pérdida de
evaluación bajó de 1,233 a 1,121, y ante la frase «La [MASK] llegó al 3\% mensual según el
INDEC» el modelo propuso «inflación» con probabilidad 0,81: ya hablaba el idioma de la prensa
económica argentina.

\textbf{Ajuste fino.} Para que la clase neutra no dominara, se tomó una muestra balanceada de
hasta 2.400 oraciones por clase (todas las de MIXTO), dividida en 80\% para entrenamiento,
10\% para validación y 10\% para prueba (6.588, 824 y 824 oraciones). Como aun así la clase
MIXTO tiene menos ejemplos que las demás, durante el entrenamiento un error en una clase escasa
se penalizó algo más que un error en una clase abundante: cuanto más rara es la clase, más pesa
cada uno de sus errores, aunque de forma moderada, para que el modelo no la ignore sin
exagerar su importancia\footnote{Técnicamente, la pérdida se ponderó por la raíz cuadrada de la
frecuencia inversa de cada clase.}. Se usaron 5 épocas, lotes de 32, tasa de aprendizaje
$2\cdot10^{-5}$ con 10\% de calentamiento, decaimiento de pesos 0,01 y selección de la mejor
época por F1 macro en validación. El F1 macro —el promedio simple del F1 de cada clase— es la
métrica principal, porque no premia acertar solo la clase más frecuente.

% ---------------------------------------------------------------------
\section{Evaluación: cuánto mide bien el termómetro}

La primera evaluación informó un F1 macro de 0,7296 en cuatro clases. Al revisarla se
encontraron tres problemas: la división de los datos se había hecho por oración y no por
noticia, de modo que las oraciones de prueba venían de noticias que también aportaban
oraciones al entrenamiento; había textos repetidos (41 oraciones de prueba, el 5\%, estaban
también en entrenamiento); y había restos del \textit{scraping} etiquetados como neutros
(«PUBLICIDAD PUBLICIDAD PUBLICIDAD», «Sobre la firma Newsletter Clarín»). Es el problema de
fuga de información descrito por \citet{kapoor2023leakage}: con noticias, la forma de dividir
los datos cambia mucho el resultado \citep{xue2026temporal}.

Para medir bien, sin modificar el modelo en uso, se limpió el conjunto etiquetado (se
descartaron 750 filas de relleno o repetidas y quedaron 10.855 oraciones), se lo dividió
\textit{por noticia} (330 noticias para entrenamiento, 39 para validación y 39 para prueba,
sin textos compartidos) y se reentrenó solo para medir, con tres semillas por división. El
código reproduce el resultado original (0,732 con la división por oración); con la división
por noticia, el F1 macro es de \textbf{0,685} (tabla~\ref{tab:finbeto}). Esa es la estimación
honesta, porque el 97\% de las noticias del corpus nunca se vio al entrenar.

Para comparar en igualdad de condiciones con el modelo de referencia,
\carpeta{finiteautomata/beto-sentiment-analysis} \citep{perez2021pysentimiento}, un BETO
ajustado con tuits en tres clases, se usaron solo las oraciones POS, NEG y NEU, contando como
error cada predicción MIXTO de FinBETO. Sobre noticias no vistas, FinBETO obtiene
\textbf{0,724} y el de referencia \textbf{0,546}: 17,7 puntos de ventaja. El de referencia
clasifica como neutras 411 de las 662 oraciones; el ajuste al dominio es lo que le permite a
FinBETO leer el tono de la prensa económica.

\begin{table}[H]
\centering\footnotesize
\begin{tabular}{@{}lcc@{}}
\toprule
\textbf{Medida (prueba, media de 3 semillas)} & \textbf{División por oración} & \textbf{División por noticia} \\
\midrule
F1 macro, 4 clases & 0,732 $\pm$ 0,012 & \textbf{0,685 $\pm$ 0,009} \\
Exactitud & 0,743 & 0,691 \\
F1 POS / NEG / NEU / MIXTO & 0,752 / 0,760 / 0,754 / 0,663 & 0,694 / 0,677 / 0,725 / 0,645 \\
\midrule
F1 macro, 3 clases: FinBETO & 0,773 & \textbf{0,724} \\
F1 macro, 3 clases: modelo de referencia & 0,620 & 0,546 \\
\bottomrule
\end{tabular}
\caption{Calidad de FinBETO según cómo se dividen los datos. La estimación que corresponde
informar es la de la división por noticia. Fuente: elaboración propia a partir de
\carpeta{Lado Noticias/sentimiento_finbert}.}
\label{tab:finbeto}
\end{table}

% ---------------------------------------------------------------------

\section{Del modelo a la temperatura de cada rueda}
\label{sec:tono-rueda}

El modelo se aplicó a todas las oraciones del corpus de entre 30 y 600 caracteres, incluido el
título. Se excluyeron 2.438 oraciones de relleno (firmas de Clarín, marcas de publicidad de
Infobae) y quedaron 399.803: 36,7\% neutras, 27,1\% positivas, 24,1\% negativas y 12,1\%
mixtas. Para cada nota se calculó el \textit{sentimiento neto}: el porcentaje de oraciones
positivas menos el de negativas; las mixtas y las neutras no suman ni restan. Con cada nota ya
asignada a su rueda (sección~\ref{sec:alineacion}) se armaron cuatro medidas:

\begin{itemize}
  \item \textbf{Tono:} el promedio del sentimiento neto de las notas de la rueda, también en una
        versión \textit{ajustada por portal} (correlación 0,995 con la original).
  \item \textbf{Tono anormal:} el tono de la rueda menos el promedio de las 20 ruedas previas.
        Mide la sorpresa en lugar del clima de fondo, y es la variable de tono de los modelos.
  \item \textbf{Atención:} la cantidad de notas de la rueda comparada con lo habitual, primero
        con las ruedas que cubren los mismos días (un lunes junta el fin de semana) y después
        con las 20 previas. En la \textit{atención anormal}, 0 es lo habitual, $+0{,}69$ el doble
        de notas y $-0{,}69$ la mitad.
  \item \textbf{Ruedas de tono extremo:} se compararon dos familias de reglas contra el azar,
        mezclando las notas entre las ruedas. La regla por desvío de los días previos marca casi
        lo mismo que el azar (15 shocks contra 16,1 con umbral 2). La regla \textit{por
        muestreo}, que tiene en cuenta cuántas notas tiene cada rueda, marca con umbral 2,5
        \textbf{19 ruedas} (10 positivas y 9 negativas) contra 3,5 del azar, y fue la elegida.
        La regla de la serie diaria original incluía al propio día en su ventana, lo que acota
        el desvío máximo a 2,27 \citep{shiffler1988}; por eso la unión la redefinió.
\end{itemize}

El termómetro tiene una lectura clara (figura~\ref{fig:u2}). El tono medio de las 244 ruedas
fue de 3,9 puntos, con un desvío de 7,0: las noticias económicas son, en promedio, levemente
positivas. La rueda más negativa fue la del 15/04/2026 ($-13{,}5$) y la más positiva, la del
14/11/2025 ($+22{,}1$). El tono de una rueda se parece al de la anterior: la correlación entre el
tono de cada rueda y el de la rueda previa es de 0,47. Si un día las noticias son más positivas
que lo habitual, al día siguiente tienden a seguir siéndolo, aunque en menor medida: el clima de
las noticias cambia de a poco, porque los temas se sostienen varios días en los medios.

\begin{figure}[H]
\centering
\includegraphics[width=0.74\linewidth]{u2_serie.png}
\caption{Arriba: tono de cada rueda (ajustado por portal), su media de las 20 ruedas previas y
las ruedas de tono extremo (verde, positivas; rojo, negativas); las líneas punteadas son los
días electorales. Abajo: atención de los medios, es decir, cuántas notas recibió cada rueda
respecto de lo habitual. Las barras naranjas son ruedas con más notas que lo habitual y las
violetas, con menos; el eje derecho traduce la escala a veces lo habitual ($\times 1{,}5$ es un
50\% más de notas, $\times 0{,}5$ la mitad). Las primeras 20 ruedas no tienen referencia previa.
Fuente: elaboración propia (\carpeta{Unión/02_sentimiento_por_rueda}).}
\label{fig:u2}
\end{figure}

% ---------------------------------------------------------------------
\section{Reconocimiento de entidades: el tono de cada empresa}
\label{sec:ner-modelo}

El termómetro mide el tono de las noticias en general. Para saber de qué empresa habla cada
oración se usó el reconocimiento de entidades nombradas (NER), que identifica nombres de
personas, organizaciones y lugares en un texto \citep{tjongkimsang2002conll}.

\textbf{Un NER híbrido.} Se combinó el modelo \texttt{es\_core\_news\_lg} de spaCy
\citep{honnibal2020spacy} con dos \textit{gazetteers} (diccionarios curados): 160 indicadores
económicos en seis categorías y las 47 empresas del panel, \textit{canonicalizadas}: cada alias
se resuelve al ticker, para que «Banco Galicia», «Grupo Financiero Galicia» y «GGAL» cuenten como
una sola entidad. Los alias en mayúsculas se buscan respetando mayúsculas (sin esa regla, AGRO
coincidía con «agro» en 920 artículos) y se excluyeron los que son palabras comunes
(«Hipotecario», «Molinos»). La limpieza quitó el relleno del \textit{scraping}, las
autorreferencias de cada medio y los fragmentos sin valor que detecta el modelo
(tabla~\ref{tab:ner}).

\begin{table}[H]
\centering\footnotesize
\begin{tabular}{@{}lll@{}}
\toprule
 & \textbf{Antes del filtrado} & \textbf{Después} \\
\midrule
Entidades únicas & 65.352 & 17.894 \\
«PUBLICIDAD» como entidad & 2.258 menciones (primera del ranking) & 0 \\
Empresas del panel & 20, fragmentadas en variantes & 47, canonicalizadas por ticker \\
Falsos positivos del alias AGRO & 920 artículos & 0 \\
\bottomrule
\end{tabular}
\caption{Efecto de la limpieza del reconocimiento de entidades. Fuente: elaboración propia a
partir de \carpeta{Lado Noticias/ner_entidades}.}
\label{tab:ner}
\end{table}

\textbf{Cruce con el sentimiento.} Cada mención se unió con el sentimiento FinBETO de su oración,
exigiendo que el texto coincida: se cruzaron el 99,37\% de las oraciones con entidades del corpus
(243.265 de 244.809). Antes de agregar, 24 variantes de indicadores se unificaron («PIB» y
«producto bruto interno» en «PBI»). YPF es la empresa más mencionada (237 de las 244 ruedas),
seguida por BBVA, Supervielle, Edenor y Galicia; tres empresas del panel no tienen ninguna
mención (CAPX, FIPL y LONG). El sentimiento de las oraciones que nombran a una empresa es el tono
del contexto en que aparece, no una opinión sobre ella.

\textbf{El tono propio.} El tono de una empresa en una rueda es (oraciones positivas $-$
negativas) / oraciones que la mencionan $\times$ 100. Como la mediana es de dos oraciones por
rueda, y en el 37\% de las ruedas hay una sola, ese valor sale muchas veces todo o nada. Por eso
se lo \textit{encoge} hacia el tono del mercado según la cantidad de oraciones $n$
\citep{efron1975}:
\begin{equation}
\text{tono propio}_{i,t} = \frac{n_{i,t}}{n_{i,t} + 10{,}9}\,\bigl(\text{tono}_{i,t} - \text{tono}_{M,t}\bigr),
\end{equation}
de modo que con una oración pesa un 8\%, con diez un 48\% y con treinta un 73\%. El tono propio
mide cuánto mejor o peor se habla de la empresa que del mercado, vale 0 si no hay noticias de
ella y se calcula para las 17 empresas con noticias frecuentes (figura~\ref{fig:u3}).

\begin{figure}[H]
\centering
\includegraphics[width=0.55\linewidth]{u3_cobertura.png}
\caption{Arriba: ruedas con al menos una noticia de cada empresa (en azul, las 17 con tono
propio). Abajo: tono propio de esas 17 empresas en cada rueda. Fuente: elaboración propia
(\carpeta{Unión/03_sentimiento_por_empresa}).}
\label{fig:u3}
\end{figure}
